\documentclass{momentum}

\title{Recognition Across Independent Authorities}
\author{Raeez Lorgat}
\date{September 2026}
\shorttitle{Recognition Across Authorities}

\begin{document}
\maketitle

\begin{abstract}
A registry can reuse a foreign identity record while retaining its own licensing decision.
Reliable reuse requires the source assertion, receiving rule, and current authority to remain separately identifiable.
This article explains how recognition preserves those records through intermediaries, policy changes, and unfinished work.
It also examines shared reservations and disclosures whose combined history reveals more than their individual answers.
A finite policy criterion identifies when knowledge, privacy, permission, and success can hold together.
Accession can reproduce bilateral relations under explicit assumptions about common terms, notice, refusals, and willingness.

The underlying results concern stated mathematical models.
Their use requires authentic instruments, faithful interpretation, complete observations, and the relevant service and provider evidence.
\end{abstract}

\section{A foreign statement and a local license}
\label{sec:registry}

North has checked a company's identity and standing.
The company applies for a license in South, which permits use of North's identity evidence.
South also requires a current licensing decision under its own rules.
Importing the file answers the identity question while the licensing question remains pending.
An authorized South decision can then permit the license, using the imported evidence.
The licensing act creates the reporting duty specified by South's rule.

This is a constructed case from \emph{The Sovereign Jurisdiction Network}, Example~4.5~\cite{sjn}.
It contains three distinct events: importing evidence, deciding the application, and issuing the license.
Repeating the import creates neither a second source assertion nor another reporting duty for the same trigger.
If South's current licensing decision refuses the application, the imported identity record leaves that refusal effective.
Proposition~4.2 proves this separation for the paper's recognition construction.

\emph{Recognition} means a receiving authority gives specified treatment or effect to another authority's determination under its own governing rule.
\emph{Admission} is the check that permits a proposed use or change under that rule.
\emph{Authority} is permission to perform the particular institutional act within its stated scope and time.
A signature identifies an attributable assertion under a key.
Its institutional effect also depends on the signer's authority and the instrument governing the requested use.

A \emph{corridor} records a directed recognition rule.
It names the source, destination, permitted evidence, treatment, and questions that require a receiving authority's current answer.
Those questions are called \emph{fresh questions}.
Fresh means that the answer must satisfy the present request and rule.
A standing delegation can authorize its computation from reusable evidence.

Domains covered by recognition can also contain fresh questions.
In the example, identity evidence and the local licensing question concern one domain.
Classifying that domain as recognized cannot remove the local question.
The complete record must retain both fields, together with the authority responsible for the answer.
\emph{How Compliance Composes}, Section~5.6, establishes this preservation through complete corridor contracts~\cite{compliance}.

South can therefore reuse North's examination while retaining responsibility for its own rule, decision, and resulting duties.
The recognition record preserves evidence of this division of work throughout later use.

\section{Keeping a decision reproducible}
\label{sec:history}

Suppose a statement was issued at time~2 and permits use through time~5.
At time~7, its retained signature can still authenticate the earlier assertion.
The expired interval prevents a use that requires current validity.
A new receipt preserves the original interval.
It cannot renew the source assertion merely by recording another delivery~\cite[Section~2.2, Proposition~2.5]{sjn}.

The verifier therefore returns separate answers about retained bytes, historical verification, current reliance, and operational authorization.
Retained bytes are the actual documents and signed records needed for verification.
A digest identifies their contents, but cannot replace missing contents.
Historical verification evaluates the original act against its recorded evidence view.
Current reliance evaluates the proposed destination, purpose, time, and policy.
Operational authorization checks the further local act and the stage at which authority must hold.

Missing evidence leaves its dependent check pending with a named unresolved premise.
A known expired interval supplies a reason for refusal of that use.
Both outcomes preserve the earlier assertion.
An \emph{evidence cutoff} records which facts were available to the historical assessment.
Effective time records when a fact governs an act.
Keeping these times separate permits a later correction without silently changing the earlier evidence view.

The source paper separates issuance, acceptance, local decisions, rights, and cash records.
Each authorized event appends to its specified target.
Earlier entries remain available when a correction, withdrawal, or reversal occurs.
The preservation theorem proves attribution and reproducible history for every finite sequence admitted under its six conditions~\cite[Section~3, Theorem~3.1]{sjn}.
Those conditions include valid authority, exact signed contents, current guards, protected dependencies, and immutable ordered history.

A \emph{guard} contains the conditions a proposed event must satisfy.
Checking the destination record alone is insufficient when its decision also depends on another authority's policy.
That policy might change after evaluation and before commitment.
The protected operation must include every changing dependency and check authority at the required stage.
A permission consumed at one stage retains only the effect its governing instrument assigns to that consumption.
Later dispatch or provider acceptance can require another current-authority check.

These are conditional guarantees about admitted records.
An unreported legal change cannot become known through an unchanged local record.
The authority source must supply the currentness evidence required by the governing use contract.

\section{What survives an intermediary}
\label{sec:routes}

A third authority may receive North's statement through South.
Its route retains the original assertion and the ordered recognition rules used along the way.
Each receiving authority supplies its own treatment and answers its own reserved questions.
Routing the same assertion through two intermediaries still supplies one distinct source.

A \emph{grade} summarizes evidence treatment on a declared ordered scale.
A treatment map describes how a receiving rule preserves or lowers an incoming grade.
The papers require compatible meanings on the common scale.
A missing correspondence between local meanings leaves transport unresolved.
Labels such as Full, Partial, and Conditional describe specified maps and conditions.
Their names alone establish no universal order between executable policies.

Consider the successful witness in \emph{Accession Networks}, Section~6.3~\cite{accession}.
The original source supplies full support.
An intermediary carries it under a condition whose admitted answer is true.
The destination applies a partial cap, and the filing requires that capped support.
The retained assertion, both policies, condition answer, and current-use certificates jointly support the filing.
Each receiving authority supplies its required answers, and the destination records Allow.

A false condition supplies no carried support.
A destination refusal prevents execution despite the same original source history.

The restriction calculation and the evidence calculation remain distinct.
In the paper's grade model, repeating an evaluation changes nothing, improving the inputs cannot worsen the result, and every participating restriction remains binding.
These conditions force combination by the greatest common lower bound, called the \emph{meet}~\cite[Section~2.1, Proposition~2.2]{compliance}.
On one ordered chain, this is the minimum.

Complementary admissible evidence can instead contribute to a supported summary through its least common upper bound, called the \emph{join}.
That summary requires a compatible witness family for the exact request.
Selecting favorable coordinates from incompatible alternatives does not provide such a family~\cite[Sections~3.3 and~5.5]{compliance}.

Applicability also remains attributable.
If one authority finds a violation and another finds its rule outside scope, the combined record retains both determinations.
The second finding does not remove the first authority's restriction.
\emph{How Compliance Composes}, Section~2.3, retains each individual source record before calculating its summary.
Repeating an identical source record changes neither its identity nor that summary.

A complete route also retains reports, deadlines, disclosure requirements, and continuing duties.
A continuing duty becomes due through its specified trigger and reaches completion through its own completion rule.
A permission result describes permitted effects.
The corresponding admitted event establishes whether an effect occurred and incurred its duties.

Replacing a route therefore requires more than equal final grades.
Two routes that both return Allow can differ because one requires another report or an earlier deadline.
The finite route comparison tests all reachable pairs of states and all declared inputs~\cite[Section~5.7, Theorem~5.25]{compliance}.
It either proves equal observation sequences or returns a distinguishing sequence.
Its conclusion assumes finite, deterministic, total machines with the same input and observation contract.
An actual evaluator needs a separate argument connecting its behavior to that model.

\section{Resources already promised elsewhere}
\label{sec:resources}

An admitted act may depend on resources that another request also needs.
A \emph{reservation} identifies capacity already committed while an operation remains unresolved.
The same physical resource needs one canonical identity across its representations.
Another receipt for that resource creates no additional capacity.

The guard construction includes every resource owner whose record can affect admission~\cite[Section~5.1]{sjn}.
It also includes the records that determine ownership and the required participants.
A changed ownership rule invalidates the earlier participant discovery.
The attempt must resolve with a terminal abort before the program recalculates dependencies and prepares again.
An unknown terminal outcome retains its reservations.
A local timeout alone cannot establish that a prepared operation will never complete elsewhere.

The capacity limit can depend on when qualified workers are available.
\emph{Accession Networks}, Example~7.15, has four inspections and one filing, with five available worker-slots~\cite{accession}.
A worker-slot is one qualified worker's reserved service during one specified interval.

The inspections can use only three slots before their deadline.
The filing can use either of two later slots.
Although total demand equals total capacity, the four inspections lack one eligible slot.
Purchasing another late filing slot leaves that deficit unchanged.
One additional qualified inspection slot before the deadline completes the assignment under the example's service assumptions.

Proposition~7.12 tests every subset of tasks against its eligible slots.
Total staffing alone cannot establish that all tasks can be assigned distinct eligible slots before their deadlines.
Existing commitments must enter the assignment before proposed additional work.
Input delivery, current authority, qualifications, and task duration remain premises of the schedule.
Each modeled task occupies one whole slot.
Its duration, including recording and handoff, must fit that slot.

Money has the same identification requirement and a separate spending restriction.
The four-inspection plan in Example~7.14 reserves twelve units of collectible collateral and separately funds four units of inspection work.
Spending the inspection fees from that collateral leaves eight pledged units and fails the stated gain bounds.
These are stipulated quantities in a finite example.
Proposition~7.13 combines collateral, service funding, and qualified time while keeping all three constraints explicit.
Its inspection and collection assumptions require their own evidence.

Recognition can make a required certificate available.
The outstanding task, reserved resource, and permitted payer still determine whether the next act can complete.

\section{What the recipient can infer}
\label{sec:privacy}

Evidence reuse can also reveal information about the source record. Selective disclosure permits a recipient to learn a stated result without receiving the complete source record.
Its privacy claim must account for what that recipient already knows and can combine with the new result.
A valid signature authenticates the disclosed answer without restricting subsequent inference.

Take three rational values $A,B,C$.
The first releases disclose $A+B$ and $B+C$.
For any rational change $t$, replacing the values by $A+t,B-t,C+t$ preserves both answers.
Each original coordinate therefore remains undetermined in this unrestricted rational model.

A further answer $A+C$ determines $A$ through
\[
 A=\frac{(A+B)+(A+C)-(B+C)}{2}.
\]
The other coordinates then follow from the earlier sums.
\emph{The Sovereign Jurisdiction Network}, Section~4.6, gives this example and its exact linear disclosure criterion~\cite{sjn}.

The criterion checks whether all observed linear equations together determine a protected quantity.
If $A$ and $B$ are nonnegative, learning $A+B=0$ already determines both.
That inference uses a domain constraint absent from the unrestricted rational test.
Separate stores must also share an admission authority when recipients can combine their answers.
Two requests can each pass against one old history and fail together.
The admitting store records each release durably before returning its answer.

The disclosure history must also account for changes to the source values.
Suppose a public operation doubles $A$ after an initial release of $A+B$.
The next sum is $2A+B$.
Subtracting the earlier answer reveals the original $A$.
Section~4.8 accounts for such public linear changes and offsets in one common history.
A new snapshot name cannot erase the recipient's earlier observation.

Answers are only part of what a recipient can observe.
Suppose possible protected values are $0,1,2$.
A query returns $0$ for value~0 and $1$ for values~1 or~2.
A gate answers only when the actual answer leaves two possibilities.
It then refuses exactly at value~0, so refusal identifies that value~\cite[Section~4.9]{sjn}.

For a deterministic mechanism, the complete finite test groups possible worlds by their entire visible output.
A \emph{world} specifies one possible source and environment history within the model.
An \emph{output fiber} contains the worlds that produce the same observation.
Each fiber must retain the required number of distinct protected values.
Reasons, pending states, timing, query choices, and visible effects all count when the observer can distinguish them.
Proposition~4.14 proves the criterion for the declared finite model and its complete observation history.

Permission to compute and permission to disclose remain separate.
A derived artifact retains its source obligations even when an algebraic simplification removes a coefficient.
Each use needs the applicable grants for its act, purpose, audience, time, and retention period~\cite[Section~4.10]{sjn}.
The inference test cannot supply those grants.

\section{When permission and privacy must hold together}
\label{sec:joined}

To satisfy privacy and permission together, an actor must choose using information it can actually observe.
If two possible worlds look identical to that actor, its deterministic policy must choose the same action in both.
An exact privacy requirement may demand equal visible actions in another set of world pairs.
Equality propagates along both kinds of connection.

The resulting \emph{joined component} contains worlds connected by either observation equality or required privacy equality.
Every component needs one action that is permitted and successful in every world it contains.
Proposition~4.18, the ``Joined permitted-success selector,'' proves this condition necessary and sufficient~\cite[Section~4.11]{sjn}.
If that common action set is empty, no such policy exists.

The paper gives four worlds, labeled $00,01,10,11$.
The actor observes the first bit.
Privacy requires equal actions for worlds sharing the second bit.
In that order, the permitted successful actions are
\[
 \{a,c\},\quad\{a,d\},\quad\{b,c\},\quad\{b,d\}.
\]

Information alone permits row choices $a,b$.
Privacy alone permits column choices $c,d$.
The joined connections link all four worlds, but no action belongs to all four sets.
The two separate checks therefore do not provide a joint policy.

An additional action $e$ admitted in every set would satisfy the combined requirement.
Its permission, funding, and successful completion must come from actual contracts in the model.
A name alone supplies none of these conditions.
The theorem fixes a finite world set and a complete visible action vocabulary.
Hidden choices, distributed participants, and randomized policies need their specified observation and execution models.

\section{Joining common terms}
\label{sec:accession}

Accession addresses how authorities establish recurring recognition relations.
Each authority joins a common instrument through a signed filing, called a \emph{deposit}.
It records permitted limits on the common terms and names counterparties it refuses.
A counterparty refusal is a \emph{declination}.
An accession reservation limits consent and differs from a reservation of operational resources.

Take three authorities $a,b,c$ and one domain.
Authorities $a$ and $c$ accept full terms, while $b$ accepts partial terms.
Authorities $b$ and $c$ decline each other.
The three deposits establish reciprocal relations between $a$ and $b$ at partial terms, and between $a$ and $c$ at full terms.
They establish no relation between $b$ and $c$~\cite[Section~2.2]{accession}.

Under five hypotheses, Theorem~4.5, ``Accession equivalence,'' gives the same pairwise relations and negotiated grades as unconstrained bilateral formation. In that comparison, every compatible pair can negotiate at zero cost~\cite[Section~4]{accession}.

The permitted reservation menu is finite and combines limits by meet.
Each reservation limits entitlements and obligations reciprocally.
Independently verifiable actual receipt of authenticated notice opens each participant's refusal window.
Subject to the instrument's remaining timing conditions, the relation takes effect after the required windows unless a timely declination prevents the pair.
A fixed rule resolves conflicting declarations.
Finally, each authority offers one common set of terms to acceptable counterparties, and its deposit accurately records those terms and refusals.

The willingness hypothesis needs evidence beyond the public rules.
If an authority offers different terms to different counterparties, a common reservation can weaken a relation or remove it.
Proposition~4.7 gives the exact condition for preserving bilateral content in that case.
Executable equivalence further requires the same versioned clauses, admitted witnesses, and local decisions for the same request and use-time context under the compared constructions~\cite[Section~6.3, Proposition~6.1]{accession}.

Notice and later updates still require work after the deposits.
Twenty sequential accessions require twenty deposits and 190 notices to incumbents.
The entrants also receive the 190 corresponding incumbent objects, so the two directions create 380 initial receipt tasks.
Two later updates by every participant create forty updates and 760 further receipt tasks.
These counts come from the lifecycle model in Section~4.5.

The construction removes repeated negotiation while notice, exceptions, assurance, and exit retain their own costs and deadlines.
It gives no general conclusion that all institutional work becomes linear or that every authority will participate.

\section{Technical reading map}
\label{sec:reading}

The following references identify the complete constructions behind this account.

\begin{description}[style=nextline]
\item[Evidence and present use.]
\emph{The Sovereign Jurisdiction Network}, Section~2.2, Proposition~2.5, separates historical verification from current reliance.
Section~3, Theorem~3.1, proves preservation under conditions N1--N6.
Section~5.1 states the complete-dependency, reservation, recovery, clock, and authority-source premises.

\item[Local questions and route histories.]
\emph{How Compliance Composes}, Section~2.3, Proposition~2.6, preserves attributable source intake.
Section~5.5, Proposition~5.13 and Lemma~5.14, preserves local decisions and source origin.
Section~5.6, Theorem~5.21, gives lossless contextual composition.
Section~5.7, Theorem~5.25, gives finite institutional route equivalence.

\item[Observation and policy composition.]
\emph{The Sovereign Jurisdiction Network}, Sections~4.6--4.11, contains the linear, public-evolution, complete-output, permission, and joined-policy constructions.
The key statements are Propositions~4.9, 4.13, 4.14, and~4.18.
Sections~4.11.1--4.11.2 give the separate randomized-policy and complete-history conditions.

\item[Accession and reserved capacity.]
\emph{Accession Networks}, Section~4.1, states H1--H5, and Section~4.2 proves Theorem~4.5.
Section~4.3 gives Proposition~4.7's content boundary.
Section~4.5 accounts for lifecycle work.
Section~6.3 gives Proposition~6.1's executable extension.
Section~7.8 gives Propositions~7.12--7.13 and Examples~7.14--7.15 on qualified time, service funding, and collateral.
\end{description}

The written proofs establish their stated implications.
The finite checks cover their supplied models and inputs.
Each application must establish authentic source records, faithful instrument interpretation, complete observations, durable implementation, and institutional performance.
The papers' status sections retain these boundaries: Section~10 in each of the three references.

\begin{thebibliography}{9}
\bibitem{sjn}
Raeez Lorgat.
\emph{The Sovereign Jurisdiction Network}.
September 2026.

\bibitem{compliance}
Raeez Lorgat.
\emph{How Compliance Composes}.
September 2026.

\bibitem{accession}
Raeez Lorgat.
\emph{Accession Networks}.
September 2026.
\end{thebibliography}

\end{document}
